\newglossaryentry{infosys}{
    name={Информационная система},
    description={совокупность аппаратных и программных компонентов, обеспечивающих сбор, хранение, обработку и выдачу информации для решения задач управления}
}

\newglossaryentry{rest-api}{
    name={REST API},
    description={архитектурный стиль взаимодействия распределённых компонентов сетевых приложений, основанный на протоколе HTTP и принципах REST}
}

\newglossaryentry{machinelearning}{
    name={Машинное обучение},
    description={раздел искусственного интеллекта, изучающий методы построения моделей, способных обучаться на данных и делать прогнозы или принимать решения без явного программирования}
}

\newglossaryentry{timeseries}{
    name={Временной ряд},
    description={последовательность данных, собранных или измеренных в последовательные моменты времени через равные или неравные промежутки}
}

\newglossaryentry{lactation}{
    name={Лактация},
    description={процесс образования и выделения молока молочной железой сельскохозяйственных животных}
}

\newglossaryentry{estrus}{
    name={Охота (эструс)},
    description={период половой активности у самок млекопитающих, во время которого происходит овуляция и возможно осеменение}
}

\newglossaryentry{mastitis}{
    name={Мастит},
    description={воспалительное заболевание вымени коровы, вызываемое бактериальной инфекцией и приводящее к снижению качества и объёма молока}
}

\newglossaryentry{ketosis}{
    name={Кетоз},
    description={метаболическое нарушение у коров в ранний послеродовой период, характеризующееся повышенным содержанием кетоновых тел в крови}
}

\newglossaryentry{culling}{
    name={Выбраковка},
    description={процесс исключения животного из стада по производственным, ветеринарным или экономическим причинам}
}

\newglossaryentry{z-score}{
    name={Z-оценка},
    description={стандартизованное отклонение значения от среднего: $z = (x - \mu) / \sigma$, где $\mu$ --- среднее по стаду, $\sigma$ --- стандартное отклонение. Значения $|z| > 2$ считаются статистически аномальными}
}

\newglossaryentry{icp}{
    name={Индекс здоровья},
    description={комплексный показатель состояния животного, рассчитываемый как взвешенная сумма Z-оценок надоя, жвачки, активности и SCC. Принимает значения от 0 до 100, где $\geq 80$ --- низкий риск, $< 40$ --- критический}
}

\newglossaryentry{confidence-interval}{
    name={Доверительный интервал},
    description={диапазон значений, в котором с заданной вероятностью находится истинное значение прогнозируемой величины. В данной работе используются квантильные интервалы q10/q90, содержащие 80~процентов распределения прогнозов}
}

\newglossaryentry{shap}{
    name={SHAP},
    description={SHapley Additive exPlanations --- метод объяснимости машинного обучения, основанный на значениях Шепли из теории кооперативных игр. Показывает вклад каждого признака в прогноз модели}
}

\newglossaryentry{ensemble}{
    name={Ансамбль моделей},
    description={метод объединения прогнозов нескольких моделей в один итоговый прогноз. В данной работе используется взвешенное среднее с весами, обратно пропорциональными MAPE каждой модели}
}
